%---------
%---RESUMO

\thispagestyle{empty}
\section*{\normalsize \centering \textbf{RESUMO}}

\noindent  

Este trabalho aborda a modernização dos processos de Extração, Transformação e Carga (ETL) da Secretaria de Estado da 
Fazendada Paraíba (SEFAZ-PB), motivada pelas limitações da arquitetura legada baseada em SQL Server Integration 
Services (SSIS) frente às demandas dos ambientes de Big Data. Nesse contexto, o estudo teve como objetivo propor 
uma reestruturação da estratégia de ETL da SEFAZ-PB, utilizando o projeto BDMalhas como estudo de caso. Para isso, 
foi desenvolvido um novo sistema de ETL baseado em Apache Airflow, Apache Spark e ClickHouse, cuja avaliação 
considerou métricas de manutenibilidade, desempenho e utilização de espaço em disco, além da análise dos desafios 
e limitações da reestruturação. Para tanto, a pesquisa foi conduzida como um estudo de caso, de caráter aplicado 
e abordagem quantitativa, envolvendo a implementação da arquitetura proposta e sua avaliação por meio de 
experimentos comparativos. Ao final, os resultados evidenciaram redução da complexidade estrutural dos pipelines, 
do tempo de execução das cargas e do espaço em disco utilizado, além de ganhos em manutenibilidade e escalabilidade.
Por outro lado, também foram identificados desafios, destacando-se a migração da arquitetura legada e as limitações 
do ambiente experimental. Por fim, conclui-se que a reestruturação proposta constitui uma alternativa viável para 
modernizar os processos de ETL da SEFAZ-PB, proporcionando ganhos em desempenho, manutenibilidade e escalabilidade, 
embora os resultados estejam restritos ao contexto do projeto BDMalhas e ao ambiente utilizado na pesquisa.

\vspace*{0.5cm}

\noindent\textbf{Palavras-chave:} ETL; Big Data; engenharia de dados; reestruturação de ETL; arquitetura de dados.
 % no máximo 4 palavras-chave separadas entre si por ponto 
 \newpage
